% Generated draft: requires mathematical and source review.
\documentclass[11pt]{article}
\usepackage[margin=25mm]{geometry}
\usepackage{fontspec}
\setmainfont{Latin Modern Roman}
\usepackage{amsmath,amssymb,mathtools,unicode-math}
\setmathfont{Latin Modern Math}
\usepackage{enumitem,needspace,xurl,hyperref}
\setlength{\parindent}{0pt}
\setlength{\parskip}{5pt}
\setlength{\emergencystretch}{3em}
\newcommand{\N}{\mathbb{N}}
\newcommand{\Z}{\mathbb{Z}}
\newcommand{\Q}{\mathbb{Q}}
\newcommand{\R}{\mathbb{R}}
\newcommand{\C}{\mathbb{C}}
\newcommand{\F}{\mathbb{F}}
\newcommand{\A}{\mathbb{A}}
\newcommand{\PP}{\mathbb{P}}
\newcommand{\E}{\mathbb{E}}
\newcommand{\eps}{\varepsilon}
\newcommand{\dd}{\,\mathrm{d}}
\begin{document}
\section*{UnsolvedMath ID: 600006}
\textbf{Source problem code:} catalogue entry
\subsection{Definitions and mathematical statement}
Let $n\geq1$ be an integer, and let $M\subset\R^{2n}$ be a smooth strictly convex hypersurface. The source asks for lower bounds on the number of $p$-periodic outer-billiard orbits for integer periods $p\ne3$.

To describe the geometric setting, equip $\R^{2n}$ with its standard symplectic form $\omega=\sum_{j=1}^{n}\dd q_j\wedge\dd r_j$. The characteristic direction at $m\in M$ is the one-dimensional kernel of $\omega$ restricted to $T_mM$. The usual geometric outer-billiard construction reflects a point $z$ through a tangency point $m\in M$ along this direction, so that the next point is $2m-z$. Defining a single-valued outer-billiard map $T_M$ requires the appropriate oriented tangency convention and domain. These are not specified in the supplied record. For a chosen such map, a periodic trajectory satisfies
\[
 z_{i+1}=T_M(z_i),\qquad i\in\Z/p\Z.
\]
The problem is to obtain lower bounds for the number of such orbits, for positive integer $p\ne3$, in the stated smooth strictly convex class.

The record does not specify whether $M$ must be closed and bound a compact convex body, whether ``$p$-periodic'' means least period exactly $p$ or period dividing $p$, or which identifications define distinct orbits. It also does not specify whether the intended bounds are uniform in $M$, their dependence on $n,p$, or any nondegeneracy assumption. These are material qualifications to resolve before a unique numerical counting assertion can be formulated. No stronger curvature or genericity assumption is imposed here.
\subsection{Short English statement}
For multidimensional outer billiards outside a smooth strictly convex hypersurface, find guarantees on the number of periodic trajectories for periods other than three. The supplied 2015 list presents the problem as open; the registry later labels it partially solved, with current status still requiring review.

The supplied research summary proposes a variational formulation on the cyclic configuration space
\[
 \{(x_1,\ldots,x_p)\in(S^{2n-1})^p:x_i\ne x_{i+1}\},\qquad x_{p+1}=x_1,
\]
using a generating function $h$ and $F=\sum_i h(x_i,x_{i+1})$. It suggests that cohomology and Morse--Lusternik--Schnirelmann methods could yield a bound of the shape $(p-1)(2n-2)+1$ for prime $p$, conditional on resolving the compactness issue near collision diagonals. This is a proposed route in the supplied summary, not an established lower bound in this draft. The summary is truncated mid-sentence, so its full hypotheses and conclusion are unavailable; its assertion that strict convexity identifies $M$ with a sphere via the Gauss map also requires clarification of the regularity and curvature conventions.
\subsection{Sources}
\begin{itemize}
\item UnsolvedMath permanent ID 600006, source code catalogue entry, ``Baker's Dozen --- Periodic multidimensional outer-billiard trajectories.'' Registry curation and metadata: Ulam AI and UnsolvedMath Contributors, CC BY 4.0; record updated 2026-08-06.
\item Tabachnikov, \emph{A Baker's Dozen of Problems} (2015), Problem 2, as identified by the supplied AMR-indexed list: \url{https://amj.math.stonybrook.edu/html-articles/Files-2015-2024/14-01/}. Publisher HTML extraction accessed 2026-07-29. The record marks status evidence and rights review as NEEDS\_REVIEW.
\item The supplied, truncated research summary attributes a cyclic-configuration-space cohomology computation to Farber--Tabachnikov (2002). It supplies no complete bibliographic reference or locator for that claim; no further citation is inferred.
\end{itemize}
\end{document}
